\section{Measure comparison and change of characteristic}
\label{sec:nonstandard-transfer}

We give two comparisons for nonstandard weighted orbital integrals,
for fixed root data and outside a finite set of residue characteristics.
First, we compare the normalized
global centralizer factors in the equal-characteristic argument of
\citet[Theorem~17.7]{HalleckDubeWFLDraft}.  The factors agree even though
the centralizer tori need only be isogenous.  Second, we express the
paired local orbital-integral identity as the vanishing of a
constructible function, with common parameters and compatible measures.
The transfer principle then gives the required mixed-characteristic
identity.

\subsection{The global centralizer factor under an isogeny}
\label{subsec:global-torus-factor}

A rational identification of character spaces need not identify the
corresponding tori.  In the global nonstandard comparison, the factor
that is invariant is the quotient of the centralizer constant by the
component-group order.  The following lemma includes the degree-lattice
term needed over a function field.

\begin{lemma}[Comparison of global torus factors]
\label{lem:global-torus-factor}
Let $K=\mathbb F_q(X)$, where $X$ is a smooth projective geometrically
connected curve of genus $g$, and let $\Gamma_K$ be its absolute Galois
group.  Let $T_1,T_2$ be $K$-tori with isomorphic
rational character modules as Galois representations, and let
$\ell:\operatorname{Lie}T_1\simeq\operatorname{Lie}T_2$ be a $K$-linear
isomorphism.  Choose local Haar measures $\eta_{i,v}$ compatible under
$\ell$ on the Lie algebras, and give the smooth integral compact
subgroups volume one at almost all places.  Write $\eta_i$ for the
resulting adelic measures, using counting measure on $T_i(K)$.

For a torus $T$, put $Y=X^*(T)$, $r=\operatorname{rank}Y^{\Gamma_K}$,
and
\[
 \deg_T(t)(\chi)=\sum_v [k(v):\mathbb F_q]\,
                         v\bigl(\chi(t_v)\bigr),\qquad
 I(T)=\bigl[\operatorname{Hom}(Y^{\Gamma_K},\mathbb Z):
                       \deg_T T(\mathbb A_K)\bigr].
\]
Let $T(\mathbb A_K)^1=\ker\deg_T$, and set
\[
 \Sha^1(K,T)=\ker\!\left(H^1(K,T)\longrightarrow
                                      \prod_v H^1(K_v,T)\right),
 \qquad b(T)=\bigl|\pi_0(\widehat T^{\Gamma_K})\bigr|.
\]
If
\[
 c_{\eta}(T)=
  \operatorname{vol}_{\eta}\bigl(T(K)\backslash T(\mathbb A_K)^1\bigr)
  \,|\Sha^1(K,T)|,
\]
then
\begin{equation}
 \label{eq:torus-factor-degree-index}
 \frac{c_{\eta_1}(T_1)/b(T_1)}{c_{\eta_2}(T_2)/b(T_2)}
       =\frac{I(T_1)}{I(T_2)}.
\end{equation}
In particular the two normalized factors agree if each torus splits at
a place of degree one.  The places may be different.
\end{lemma}

\begin{proof}
We will express each normalized factor as the same scalar times
$I(T_i)$.  First, the exponential sequence for the complex dual torus is
\[
 0\longrightarrow Y\longrightarrow Y\otimes_{\mathbb Z}\mathbb C
   \longrightarrow\widehat T(\mathbb C)\longrightarrow 1.
\]
The Galois action factors through a finite group.  Taking its invariants,
and averaging on the complex vector space, identifies the component
group of $\widehat T^{\Gamma_K}$ with $H^1(K,Y)$.  Indeed, the image of
the invariant vector space is precisely the identity component of the
fixed torus.  Thus
\begin{equation}
 \label{eq:torus-component-cohomology}
 b(T)=|H^1(K,Y)|.
\end{equation}

We specify the measure comparison.  Choose invariant top differential
forms $\omega_i$ over $K$ with $\ell^*\omega_2=\omega_1$.
Compatibility of the local Haar measures means
\[
 \eta_{i,v}=a_v|\omega_i|_v
\]
with the same positive constant $a_v$ for $i=1,2$.  Put $d=\dim T_i$.
Let $P_v(T,t)$ be the local Artin polynomial and
$\rho(T)=\lim_{s\to1}(s-1)^rL(T,s)$, in the conventions of
\citet[Section~6]{GeisserSuzuki2022}.  The rational Galois-module
isomorphism gives the same $d,r,P_v$, and $\rho$ for both tori.  In these
conventions the Tamagawa measure is
\begin{equation}
 \label{eq:torus-tamagawa-measure}
 \mu_i=\frac{1}{\rho\,q^{(g-1)d}}
              \prod_v P_v(q_v^{-1})^{-1}|\omega_i|_v,
 \qquad q_v=|k(v)|.
\end{equation}
Consequently $\eta_i=\alpha\mu_i$ with the same scalar
\[
 \alpha=\rho\,q^{(g-1)d}\prod_v a_vP_v(q_v^{-1}).
\]
At almost all places the forms are integral generators and the tori
have good reduction; there the volume of the integral compact under
$|\omega_i|_v$ is $P_v(q_v^{-1})$.  The displayed product therefore has
only finitely many factors different from $1$.  Replacing a global form
by a $K^\times$-multiple does not change the global measure, by the product
formula.  This also explains why clearing denominators in a rational
torus correspondence adds no unrecorded global scalar.

For clarity, the function-field Tamagawa number includes a degree
correction.  The definition in \citet[Section~6]{GeisserSuzuki2022} and
their Proposition~6.3 give
\[
 \operatorname{vol}_{\mu_i}
        \bigl(T_i(K)\backslash T_i(\mathbb A_K)^1\bigr)
   =(\log q)^r I(T_i)\tau(T_i),
 \qquad
 \tau(T_i)=\frac{|H^1(K,X^*(T_i))|}{|\Sha^1(K,T_i)|}.
\]
Here their discriminant $\operatorname{Disc}(h_{T_i})$ is exactly
$I(T_i)$: the degree pairing identifies the torsion-free adelic class
group with the image lattice of $\deg_{T_i}$ in
$\operatorname{Hom}(X^*(T_i)^{\Gamma_K},\mathbb Z)$.
Together with \eqref{eq:torus-component-cohomology}, this proves the
more explicit formula
\begin{equation}
 \label{eq:torus-factor-explicit}
 \frac{c_{\eta_i}(T_i)}{b(T_i)}
       =\alpha(\log q)^r I(T_i),
\end{equation}
and hence \eqref{eq:torus-factor-degree-index}.

Finally suppose $T$ splits at a degree-one place $v_0$.
The invariant sublattice $Y^{\Gamma_K}$ is saturated in $Y$:
if $ny$ is invariant for $n\ne0$, torsion-freeness of $Y$ makes $y$
invariant.  Therefore restriction gives a surjection
\[
 X_*(T_{K_{v_0}})=\operatorname{Hom}(Y,\mathbb Z)
   \longrightarrow\operatorname{Hom}(Y^{\Gamma_K},\mathbb Z).
\]
For a local cocharacter $\lambda$ and uniformizer $\varpi_{v_0}$,
the adele equal to $\lambda(\varpi_{v_0})$ at $v_0$ and to $1$
elsewhere has degree $\chi\mapsto\langle\chi,\lambda\rangle$.
Thus the degree map is onto and $I(T)=1$.
\end{proof}

\paragraph{Application to the global nonstandard comparison.}
In \citet[Corollary~17.11, equation~(17.2.2)]{HalleckDubeWFLDraft},
take $K=F_0$ and let $T_i$ be the generic fiber of the regular
centralizer $J_{i,a_M}$.  The rational root-ray correspondence is
equivariant for the common cameral monodromy.  It gives isomorphic
rational character modules; after clearing the fixed denominators it
gives a torus isogeny.  At the good characteristics under consideration
its differential is invertible and induces the Lie-algebra
identification used for the local measures, as in
\citet[Section~1.12, before Theorem~1.12.7]{Ngo2010FundamentalLemma}.
There is no requirement that the integral character lattices be
isomorphic.

The marked point in the global construction is rational, and the
rigidification makes each $T_i$ split there
\citep[Section~14.1 and the proof of Lemma~15.11]{HalleckDubeWFLDraft}.
Hence $I(T_1)=I(T_2)=1$.  Moreover,
\citet[Lemma~12.2.5, pp.~1735--1736]{ChaudouardLaumon2012II}
identify the Frobenius-coinvariant kernel used in that draft with
$\Sha^1(F_0,T_i)$:
\[
 \ker\bigl(T_i(F)_\tau\longrightarrow T_i(\mathbb A)_\tau\bigr)
       \simeq\Sha^1(F_0,T_i).
\]
Here $F=\overline{\mathbb F}_q(X)$, $\mathbb A$ is its geometric
adele ring in that construction, and $\tau$ is arithmetic Frobenius.
This statement concerns coinvariants for the infinite cyclic Weil
group; it does not replace them silently by continuous Galois
cohomology.  Also $\widehat T_i^{\Gamma_{F_0}}$ is the fixed dual torus
denoted $\widehat T_i^{W_a}$ in the trace formula.  To see this even
when $W_a$ is defined geometrically, use the section of the arithmetic
fundamental group supplied by the rational marked point.  Its image
acts trivially on the character lattice because the torus is split
there.  The arithmetic and geometric monodromy images on that lattice
are consequently equal.  The lemma therefore
gives the exact cancellation needed in equation~(17.2.2):
\begin{equation}
 \label{eq:nonstandard-centralizer-cancellation}
 \frac{c(J_{1,a_M})}{|\pi_0(\widehat T_1^{W_a})|}
       =\frac{c(J_{2,a_M})}{|\pi_0(\widehat T_2^{W_a})|}.
\end{equation}
The two numerators and the two denominators need not agree separately.
Equation~\eqref{eq:nonstandard-centralizer-cancellation} gives their
combined cancellation in place of the generic-torus identifications
and separate cancellations on page~169 of that draft.

\paragraph{The remaining coroot covolumes.}
It remains to identify the factors in its equation~(17.2.3) with
Waldspurger's coefficient.  This calculation uses the coroot lattices
within each group, rather than an identification of integral torus
lattices across the pair.  Let $S_i$ be the pinned maximal torus of
$G_i$, and let $\Gamma$ be the based-root-datum monodromy group.
Denote the rational root-ray correspondence by $\psi$, and write
\[
 \psi_*:X_*(S_1)\otimes_{\mathbb Z}\mathbb Q
       \xrightarrow{\ \sim\ }
       X_*(S_2)\otimes_{\mathbb Z}\mathbb Q
\]
for its $\Gamma$-equivariant map on geometric cocharacters.
For $H_i=M_i,G_i$, put
$\Lambda_{H_i}=\mathbb Z\langle\Phi_{H_i}^{\vee}\rangle$ inside the
geometric cocharacter lattice $X_*(S_i)$.
The map $\psi_*$ carries $\Lambda_{H_1}\otimes\mathbb Q$ onto
$\Lambda_{H_2}\otimes\mathbb Q$ for $H=M,G$, and therefore induces
isomorphisms on the corresponding rational coinvariant quotients.
The lattice and real space in
\citet[Lemma~17.10]{HalleckDubeWFLDraft} are
\[
 \mathcal X_*(H_i)=
       \bigl((X_*(S_i)/\Lambda_{H_i})_\Gamma\bigr)_{\mathrm{tf}},
 \qquad
 \mathfrak a_{H_i}=\mathcal X_*(H_i)\otimes_{\mathbb Z}\mathbb R,
\]
where the subscript $\mathrm{tf}$ means quotient by torsion.
For $S_i$ the same definition uses the zero coroot lattice.
Write $\mathfrak a_{S_i}^{H_i}$ for the kernel of
$\mathfrak a_{S_i}\to\mathfrak a_{H_i}$, with its induced measure.
We also write $\psi_*$ for the induced real isomorphisms
$\mathfrak a_{H_1}\simeq\mathfrak a_{H_2}$ and
$\mathfrak a_{S_1}^{H_1}\simeq\mathfrak a_{S_2}^{H_2}$.
Set
\[
 V_{H_i}=\operatorname{vol}\bigl(
  \mathfrak a_{S_i}^{H_i}/(\Lambda_{H_i})_\Gamma
                                  \bigr).
\]
First suppose $\mathfrak a_{G_i}=0$.  Dividing the formula of that
lemma for $H_i=G_i$ by the formula for $H_i=M_i$ gives
\[
 \frac{V_{G_i}}{V_{M_i}}
  =\frac{|\pi_0(Z_{\widehat G_i}^{\Gamma})|}
          {|\pi_0(Z_{\widehat M_i}^{\Gamma})|}
       \operatorname{vol}\bigl(\mathfrak a_{M_i}/\mathcal X_*(M_i)\bigr).
\]
The torus covolume and the torus component factor cancel within each
fixed $i$.  Transporting the real measures through $\psi_*$ gives
\[
 \frac{V_{H_1}}{V_{H_2}}=
 [ (\Lambda_{H_2})_\Gamma:
                  \psi_*((\Lambda_{H_1})_\Gamma)],\qquad H=M,G.
\]
Here $[A:B]=[A:A\cap B]/[B:A\cap B]$ is the relative index of
commensurable lattices: it is the covolume of $B$ divided by that of
$A$ in their common real space.  Thus the remaining ratio is
Waldspurger's coefficient
\[
 c_\psi=\frac{V_{G_1}V_{M_2}}{V_{G_2}V_{M_1}}=
 \frac{[(\Lambda_{G_2})^\Gamma:
                         \psi_*((\Lambda_{G_1})^\Gamma)]}
      {[(\Lambda_{M_2})^\Gamma:
                         \psi_*((\Lambda_{M_1})^\Gamma)]}.
\]
The simple coroots form permutation bases for
$\Gamma$, and $\psi_*$ scales their rays by scalars constant on each
$\Gamma$-orbit.  It acts with the same scalars on orbit sums in the
invariants and orbit classes in the coinvariants.  The relative
indices therefore agree, as in the proof of
\citet[Corollary~17.11]{HalleckDubeWFLDraft}.
For general $G_i$, that proof reduces to this case by quotienting by
the maximal split central tori.  Its reduction, using Theorem~15.16,
preserves the weighted orbital integrals, the stable recursions, and
the coroot-lattice coefficient.

This proves the required global centralizer normalization.  The
geometric inputs remain the isogeny statement for neutral Picard
components and the cohomological comparison on the anisotropic locus
\citep[Proposition~4.18.1 and Theorem~8.8.2]{Ngo2010FundamentalLemma},
together with the support statement used in the draft.  The split
rational marked point identifies arithmetic and geometric monodromy
\citep[Proposition~2.21 and Corollary~2.22]{HalleckDubeWFLDraft};
thus, after the split-center reduction, the center satisfies the
geometric anisotropy hypothesis of that comparison.  The calculation
does not assert a generic isomorphism $J_1^0\simeq J_2^0$ or constitute
a new proof of those geometric results.

\subsection{Change of characteristic}

We now transfer the local nonstandard identity, with the compatible
normalizations just described.  The uniformity is over a fixed finite
collection of root data and all its integral generically regular
matching parameters.

\begin{lemma}[Change of characteristic for fixed nonstandard data]
\label{lem:nonstandard-change-characteristic}
Fix a finite collection of paired based root data with rational
root-ray maps $\psi$, matching Levi data, and compatible finite cyclic
outer actions.  Suppose the nonstandard weighted fundamental lemma
holds over local fields of sufficiently large positive characteristic
for their unramified forms, including all cyclic-generator variants.
Then it holds over characteristic-zero local fields of sufficiently
large residue characteristic, uniformly for this collection and for
all integral generically regular matching parameters.

Here the ambient hyperspecial subgroups have volume one; the Arthur
weight-space measures are transported through $\psi$; and the paired
regular-centralizer measures are transported through the induced
Lie-algebra isomorphism.
\end{lemma}

\begin{proof}
We apply the general transfer principle for parameterized integrals
\citep[Theorems~2.8.1--2.8.3]{CluckersHalesLoeser2011}.
To do so, we construct a common parameter space for the paired
integrals, express their measures by differential forms, and remove
the volume and discriminant normalizations.  The resulting identity
will be the vanishing of one constructible function.

\paragraph{Common parameters and the identity.}
Choose an integer $N$ excluding the bad primes for the finitely many
root data, Chevalley and Kostant constructions, and outer actions,
and making $\psi$ and its inverse integral over $R=\mathbb Z[1/N]$.
Include the ordinary endoscopic descendants in this finite list.
Write $G_i$ and $M_i$ for the groups and Levi subgroups attached to
the two paired data, and $\mathfrak t_i$ for their chosen Cartan Lie algebras.
The Weyl-equivariant Cartan identifications of
\citet[Sections~1.12.4--1.12.6]{Ngo2010FundamentalLemma} give
\[
 C_M=\mathfrak t_1/W_{M_1}\simeq\mathfrak t_2/W_{M_2},
 \qquad
 C_G=\mathfrak t_1/W_{G_1}\simeq\mathfrak t_2/W_{G_2}
\]
over $R$, compatibly with $C_M\to C_G$ and with the common outer
action.  They identify the integral models and the generically
$G$-regular loci.  Consequently the same statements hold after a
common unramified outer twist.

For a local field $F$ and an integral generically regular parameter $b$
in this common base, let $Y_{i,b}$ be its Kostant lift to
$\operatorname{Lie}(M_i)$.
Write $S_{M_i}^{G_i}(b)$ for the stable weighted orbital integral at
$Y_{i,b}$ with test function
$1_{\operatorname{Lie}(G_i)(\mathcal O_F)}$.
With the measures in the statement, the identity to be transferred is
\begin{equation}
 \label{eq:nonstandard-local-identity}
 S_{M_1}^{G_1}(b)=c_\psi S_{M_2}^{G_2}(b),
\end{equation}
in the convention of
\citet[Definition~16.4 and Theorem~17.7]{HalleckDubeWFLDraft}.
The coroot-lattice coefficient $c_\psi$ is fixed on each cyclic-action
stratum.  We next parameterize the unramified splitting extensions
and outer twists so that both sides of
\eqref{eq:nonstandard-local-identity} are defined on the same space.

Encode an unramified splitting extension of fixed degree $d$ by a
monic polynomial with integral coefficients and irreducible separable
reduction, together with the matrix of an algebra automorphism $\tau$
of exact order $d$.  All these conditions are first-order conditions
of fixed degree, and every unramified extension has such a
presentation.  We use the general coefficient encoding of
\citet[Section~3.2]{CluckersHalesLoeser2011}; restriction to binomials
is unnecessary.  Use the same extension and generator on both sides.
The resulting common definable space is
\[
 \Lambda_F=
 \{(\mathbf a,\tau,b):
 b\in C_{M,\mathbf a,\tau}(\mathcal O_F),\quad
 b\text{ is generically }G\text{-regular}\}.
\]
The subscript denotes the descended integral model: in extension
coordinates, its points satisfy the common semilinear fixed-point
equations.  The map $\psi$ intertwines those equations.  Thus every
required matching pair occurs over $\Lambda_F$.  Varying $\tau$ only
adds the specified cyclic-generator variants.

\paragraph{Compatible measures.}
Over this regular locus, the centralizers $J_i$ in $M_i$ are the
twists of the pinned tori along the same cameral Weyl torsor.
Multiplying $\psi$ by an integer produces an equivariant torus
isogeny.  Descending it and dividing its differential by that
integer gives an algebraic isomorphism
\[
 \lambda:\operatorname{Lie}(J_1)\simeq
                 \operatorname{Lie}(J_2).
\]
Enlarge $\Lambda$ by a nonzero top covector $\omega_1$ on
$\operatorname{Lie}(J_1)$, and set
$\omega_2=(\lambda^{-1})^*\omega_1$.
Their invariant extensions give compatible Haar measures, as in
\citet[Section~1.12, before Theorem~1.12.7]{Ngo2010FundamentalLemma}.
Every field-valued parameter has such a frame.  Any other compatible
pair differs by a common positive scalar, which rescales both sides
of the desired identity by its inverse.

For every ambient group $H$ in either finite stable recursion,
also include a nonzero invariant top-form frame $\eta_H$.
On a regular stable orbit $X_b$ in $\operatorname{Lie}(H)$, the
exact sequence
\[
 0\longrightarrow\operatorname{Lie}(J_b)
 \longrightarrow\operatorname{Lie}(H)
 \xrightarrow{Z\mapsto[Z,Y]}T_YX_b\longrightarrow0
\]
gives the quotient orbital form from $\eta_H$ and $\omega_J$.
On finite smooth charts its comparison with the motivic measure
is the absolute value of a rational determinant, hence a
constructible factor $\mathbb L^{-\operatorname{ord}(r)}$.
This is the differential-form realization of quotient measures
used in \citet[Sections~7 and~8.6]{CluckersHalesLoeser2011}.
If $K_H$ is the hyperspecial subgroup, put
\[
 V_H=\int_{K_H}|\eta_H|>0.
\]
This is a constructible integral.  The normalized ambient measure is
$|\eta_H|/V_H$; below we clear the finitely many factors $V_H$.

\paragraph{Discriminants and stable recursion.}
We must also account for the discriminant normalization
$|D_H(Y)|^{1/2}$ in the weighted orbital integrals.
In either stable recursion, write $M$ for its fixed Levi.
Every ordinary endoscopic descendant has this same actual $F$-Levi.
Choose an $F$-parabolic $P_H=MU_H$ and set
\[
 R_{H,M}(Y)=\det\bigl(\operatorname{ad}(Y)
                     \mid\operatorname{Lie}(U_H)\bigr).
\]
Opposite roots give
\[
 D_H(Y)=(-1)^{\dim U_H}D_M(Y)R_{H,M}(Y)^2.
\]
Thus division of the specialized stable recursion by
$|D_M(Y)|^{1/2}$ replaces every discriminant factor by
$|R_{H,M}(Y)|$, an integral valuation power of the residue
cardinality.  Under the pairing, the two Levi discriminants differ
by a fixed rational scalar $u_\psi$.  Enlarge $N$ by its numerator
and denominator.  Then $|u_\psi|=1$, so the two removed
Levi-discriminant factors agree at every matched parameter.

The ordinary recursion is
\[
 S_M^H=J_{M,M}^H-
   \sum_{H'\ne H}\iota_M(H,H')S_M^{H'}.
\]
Its index sets are finite and its rational coefficients depend
only on the fixed root data and cyclic action; each $H'$ has fewer
roots and the same Levi $M$
\citep[Section~9.1]{CluckersHalesLoeser2011}.
The weighted integrands and hyperspecial-lattice indicators are
the definable, relatively integrable families constructed in
\citet[Sections~6 and~8.6]{CluckersHalesLoeser2011}.
Use their integer-valuation weights with the fixed compatible
weight-space volumes; real constants are permitted by Section~2.10
of that reference.  The form comparisons and resultant factors
above preserve constructibility and relative integrability.
Indeed, the intrinsic quotient form and the Lie-algebra/Chevalley
quotient form are invariant on each geometrically transitive
regular stable orbit.  Their ratio comes from the base, including
the frames, and can be evaluated at the Kostant lift.  Its absolute
value is a constructible valuation power, so the projection formula
preserves the relative integrability of the CHL families.

Define the modified functions by this recursion using
$|R_{H,M}|$ times the quotient orbital integrals.  Induction shows
that their specializations equal $S_M^H/|D_M|^{1/2}$.
No division by a square root in the motivic ring is being made.
With this common discriminant factor removed, expand the finite
recursions and multiply \eqref{eq:nonstandard-local-identity} by the
product of the ambient volumes $V_H$, with multiplicities if necessary.
The result is the vanishing of one
constructible function $\Phi$ on the enlarged common parameter
space.  Each specialized multiplier is positive.

\paragraph{Transfer and recovery of the normalized identity.}
By hypothesis $\Phi_{k((t))}$ vanishes identically for every finite
residue field $k$ of sufficiently large characteristic, including
all permitted generator variants.  The abstract transfer principle
therefore makes $\Phi_F$ identically zero for every sufficiently
large-residue-characteristic local field $F$ with residue field
$k$, in particular every characteristic-zero such $F$.
Undo the positive ambient-volume multiplier and the common
nonzero Levi-discriminant factor.  The surjectivity of the
extension and frame parameterizations yields the required identity
for every matching parameter.  Taking the maximum of the finitely
many bounds proves the assertion.
\end{proof}
